\chapter{Verification conditions}
\source{Sample Q3, p. 3; DafnyVerificationConditions PDF, physical pp. 3--30.}
These are written partial-correctness questions, not requests to run a verifier.
State the induced formulas even if a proposed invariant is invalid. Distinguish
deriving an obligation from proving it. Ignore termination and runtime-safety
obligations in this practice set unless explicitly asked.

\task{6.1. Assignment and sequence}
Compute each weakest precondition. Show substitution, working backward.
\begin{enumerate}[label=(\alph*)]
\item \code{x := x + 3;} with postcondition $x\geq 10$.
\answer{[Substitution and simplified precondition.]}
\item \code{x := y + 2; y := x - 1;} with postcondition $x>y\land y\geq 0$.
\answer{[Precondition before each assignment; simplified result.]}
\item \code{x := x + y; y := x - y;} with postcondition $y=7$.
\answer{[Show both substitutions. Assignments are sequential.]}
\end{enumerate}

\task{6.2. Conditional}
Compute a weakest precondition for this code with postcondition $z\geq 0$.
First write the guarded implications for both branches; then simplify if useful.
\par\noindent\begin{minipage}{\linewidth}
\begin{lstlisting}
if x < y {
  z := y - x;
} else {
  z := x - y;
}
\end{lstlisting}
\end{minipage}\par
\answer{[Branch preconditions; guarded conjunction; simplification.]}

\task{6.3. Derive all three loop obligations}
\par\noindent\begin{minipage}{\linewidth}
\begin{lstlisting}
method TwicePlus(n:int) returns (t:int)
  requires n > 0
  ensures t == 2*n + 3
{
  t := 0;
  var i := n;
  while i > 0
  {
    t := t + 2;
    i := i - 1;
  }
  t := t + 3;
}
\end{lstlisting}
\end{minipage}\par
For \textbf{each} proposed invariant below, list (1) initialization after
substituting the initialization, (2) preservation after substituting the body,
and (3) exit after accounting for the final assignment. Then state which
obligations hold. For a failed implication, give a concrete counterexample
satisfying its antecedent. A counterexample need not be a reachable loop state.

\begin{enumerate}[label=(\alph*)]
\item $i\geq 0\land t+2i=2n$.
\answer{Initialization: [ ]. Preservation: [ ]. Exit: [ ]. Validity/counterexamples: [ ].}
\item $i\geq 0\land t=2i$.
\answer{Initialization: [ ]. Preservation: [ ]. Exit: [ ]. Validity/counterexamples: [ ].}
\item $i\geq 0$.
\answer{Initialization: [ ]. Preservation: [ ]. Exit: [ ]. Validity/counterexamples: [ ].}
\end{enumerate}

\task{6.4. Repair and transfer}
Now change the body update to \code{t := t + 3} and the postcondition to
$t=3n+3$, retaining the final \code{t := t + 3}. Propose an invariant and derive
the three obligations. Include bounds you need at loop exit.
\answer{Invariant: [ ]. Initialization: [ ]. Preservation: [ ]. Exit: [ ].}
